% !PRINTED 280
\weberpage{280}
%<<CONT
worin die $q_1, q_2 \ldots q_{m-1}$ und ebenso die $f_1, f_2 \ldots f_m$ ganze rationale Functionen von $x$ sind; die Grade der Functionen $f, f_1, f_2, \ldots f_m$ \gesperrt{nehmen ab} und man kann daher die Operation so weit fortsetzen, dass $f_m$ constant ist, oder wenigstens in dem betrachteten Intervall nicht mehr verschwindet. Dass $f_m$ nicht Null werden kann, ist eine Folge der Voraussetzung, dass $f$ und $f_1$ ohne gemeinsamen Theiler sind.

Dass man dann in der Reihe
\weq{3}{f,\, f_1,\, f_2 \ldots f_m}
wirklich eine Sturm'sche Kette hat, ergiebt sich unmittelbar, wenn man die Kriterien \S.\,86, 1. bis 4. durchgeht. Denn wenn zwei auf einander folgende der Functionen (3) zugleich verschwinden, so verschwinden nach (2) auch alle nachfolgenden, dies ist aber unm\"oglich, weil $f_m$ von Null verschieden ist.

Ist aber $f_\nu(x) = 0$, so folgt aus (2)
\[ f_{\nu-1}(x) = -f_{\nu+1}(x), \]
womit alle die Forderungen des \S.\,86 befriedigt sind.

Es ist, wie sich von selbst versteht, gestattet, die Functionen der Reihe (3) mit positiven, z.\,B. constanten Factoren zu multipliciren, ohne dass sie aufh\"oren, eine \gesperrt{Sturm}'sche Kette zu bilden.

F\"ur $n=2$ k\"onnen wir demnach als die Sturm'schen Functionen folgende nehmen:
\[
\begin{gathered}
f(x) = x^2 + a x + b\\
f_1(x) = 2x + a\\
f_2(x) = a^2 - 4b.
\end{gathered}
\]

\begin{center}
\S.\,90.

\gesperrt{Hermite's L\"osung des Sturm'schen Problems.}
\end{center}

Ein anderer Weg zur L\"osung des Sturm'schen Problems, der zu einfacheren Resultaten f\"uhrt, wenigstens was die Durchf\"uhrung der Rechnung im Einzelnen betrifft, ist von \gesperrt{Hermite} eingeschlagen, der an das Tr\"agheitsgesetz der quadratischen Formen und die Tschirnhausen-Transformation ankn\"upft\footnote{\gesperrt{Hermite}, Remarques sur le th\'eor\`eme de M. Sturm. Comptes rendus der Pariser Akademie, T.\,36 (1853).}.
